% ============== 第六节 财务会计信息与管理层分析 ==============================
% 审计意见 → 合并三表主要数据 → 非经常性损益 → 主要财务指标 → 管理层讨论
% （资产/偿债/经营成果/期间费用/现金流）→ 税收优惠 → 股利分配 → 盈利预测。
% 全部数据内部勾稽：资产=负债+权益、权益滚动=期初+净利润（报告期未分红）、
% 各比率可由报表数据复算（加权 ROE、EPS 按证监会编报规则第 9 号口径）。
\gssection{财务会计信息与管理层分析}

本节财务数据除特别说明外，均引自经示例会计师审计的公司财务报告。示例会计师
对公司2023年12月31日、2024年12月31日、2025年12月31日的合并及母公司资产负债
表，2023年度、2024年度、2025年度的合并及母公司利润表、现金流量表、股东权益
变动表以及财务报表附注出具了标准无保留意见的《审计报告》（示例会所审字
〔2026〕第××号）。

\subsection{合并财务报表主要数据}

\subsubsection{合并资产负债表主要数据}

\biaodw{万元}
\begin{gstab}
\begin{longtable}{|L{40mm}|R{32mm}|R{32mm}|R{32mm}|}
\hline
\thA{项目} & \thBw{32mm}{2025年12月31日} & \thBw{32mm}{2024年12月31日} & \thBw{32mm}{2023年12月31日} \\ \hline
\endfirsthead
\hline
\thA{项目} & \thBw{32mm}{2025年12月31日} & \thBw{32mm}{2024年12月31日} & \thBw{32mm}{2023年12月31日} \\ \hline
\endhead
流动资产合计 & 102,846.53 & 78,634.92 & 58,462.35 \\ \hline
非流动资产合计 & 65,545.92 & 49,840.44 & 37,922.17 \\ \hline
{\bfseries 资产总计} & {\bfseries 168,392.45} & {\bfseries 128,475.36} & {\bfseries 96,384.52} \\ \hline
流动负债合计 & 55,782.64 & 42,158.36 & 30,236.45 \\ \hline
非流动负债合计 & 15,554.54 & 11,238.08 & 7,911.66 \\ \hline
{\bfseries 负债合计} & {\bfseries 71,337.18} & {\bfseries 53,396.44} & {\bfseries 38,148.11} \\ \hline
归属于母公司股东的权益 & 97,055.27 & 75,078.92 & 58,236.41 \\ \hline
{\bfseries 股东权益合计} & {\bfseries 97,055.27} & {\bfseries 75,078.92} & {\bfseries 58,236.41} \\ \hline
\end{longtable}
\end{gstab}

\subsubsection{合并利润表主要数据}

\biaodw{万元}
\begin{gstab}
\begin{longtable}{|L{46mm}|R{30mm}|R{30mm}|R{30mm}|}
\hline
\thA{项目} & \thB{2025年度} & \thB{2024年度} & \thB{2023年度} \\ \hline
\endfirsthead
\hline
\thA{项目} & \thB{2025年度} & \thB{2024年度} & \thB{2023年度} \\ \hline
\endhead
营业收入 & 132,594.18 & 105,238.47 & 82,456.32 \\ \hline
营业成本 & 74,231.28 & 59,796.79 & 47,536.85 \\ \hline
利润总额 & 25,297.55 & 19,382.94 & 13,853.26 \\ \hline
所得税费用 & 3,321.20 & 2,540.43 & 1,817.98 \\ \hline
{\bfseries 净利润} & {\bfseries 21,976.35} & {\bfseries 16,842.51} & {\bfseries 12,035.28} \\ \hline
归属于母公司股东的净利润 & 21,976.35 & 16,842.51 & 12,035.28 \\ \hline
扣除非经常性损益后归属于母公司股东的净利润 & 21,235.68 & 16,238.47 & 11,562.83 \\ \hline
\end{longtable}
\end{gstab}

\subsubsection{合并现金流量表主要数据}

\biaodw{万元}
\begin{gstab}
\begin{longtable}{|L{46mm}|R{30mm}|R{30mm}|R{30mm}|}
\hline
\thA{项目} & \thB{2025年度} & \thB{2024年度} & \thB{2023年度} \\ \hline
\endfirsthead
\hline
\thA{项目} & \thB{2025年度} & \thB{2024年度} & \thB{2023年度} \\ \hline
\endhead
经营活动产生的现金流量净额 & 19,435.82 & 14,528.36 & 10,236.45 \\ \hline
投资活动产生的现金流量净额 & -15,246.38 & -12,835.47 & -8,452.36 \\ \hline
筹资活动产生的现金流量净额 & 4,265.00 & 3,825.00 & 2,150.00 \\ \hline
现金及现金等价物净增加额 & 8,454.44 & 5,517.89 & 3,934.09 \\ \hline
期末现金及现金等价物余额 & 26,356.89 & 17,902.45 & 12,384.56 \\ \hline
\end{longtable}
\end{gstab}

\subsection{非经常性损益明细}

报告期各期，公司非经常性损益构成如下：

\biaodw{万元}
\begin{gstab}
\begin{longtable}{|L{46mm}|R{30mm}|R{30mm}|R{30mm}|}
\hline
\thA{项目} & \thB{2025年度} & \thB{2024年度} & \thB{2023年度} \\ \hline
\endfirsthead
\hline
\thA{项目} & \thB{2025年度} & \thB{2024年度} & \thB{2023年度} \\ \hline
\endhead
计入当期损益的政府补助 & 862.45 & 702.35 & 545.60 \\ \hline
委托他人投资或管理资产的损益 & 45.82 & 38.26 & 32.15 \\ \hline
其他营业外收入和支出 & -35.28 & -28.64 & -21.45 \\ \hline
小计 & 872.99 & 711.97 & 556.30 \\ \hline
减：所得税影响额 & 132.32 & 107.93 & 83.85 \\ \hline
{\bfseries 非经常性损益净额} & {\bfseries 740.67} & {\bfseries 604.04} & {\bfseries 472.45} \\ \hline
\end{longtable}
\end{gstab}

\subsection{主要财务指标}

\begin{gstab}
\begin{longtable}{|L{40mm}|R{32mm}|R{32mm}|R{32mm}|}
\hline
\thA{项目} & \thBw{32mm}{2025年度／\newline 2025年12月31日} & \thBw{32mm}{2024年度／\newline 2024年12月31日} & \thBw{32mm}{2023年度／\newline 2023年12月31日} \\ \hline
\endfirsthead
\hline
\thA{项目} & \thBw{32mm}{2025年度／\newline 2025年12月31日} & \thBw{32mm}{2024年度／\newline 2024年12月31日} & \thBw{32mm}{2023年度／\newline 2023年12月31日} \\ \hline
\endhead
流动比率（倍） & 1.84 & 1.87 & 1.93 \\ \hline
速动比率（倍） & 1.40 & 1.39 & 1.42 \\ \hline
资产负债率（合并） & 42.36\% & 41.56\% & 39.58\% \\ \hline
应收账款周转率（次） & 4.91 & 5.00 & 4.97 \\ \hline
存货周转率（次） & 3.34 & 3.39 & 3.34 \\ \hline
研发投入占营业收入的比例 & 9.72\% & 9.34\% & 8.99\% \\ \hline
基本每股收益（元／股） & 0.61 & 0.47 & 0.33 \\ \hline
稀释每股收益（元／股） & 0.61 & 0.47 & 0.33 \\ \hline
加权平均净资产收益率 & 25.53\% & 25.27\% & 23.05\% \\ \hline
扣除非经常性损益后加权平均净资产收益率 & 24.67\% & 24.36\% & 22.14\% \\ \hline
\end{longtable}
\end{gstab}
\tabnote{每股收益及加权平均净资产收益率按照中国证监会《公开发行证券的公司信
息披露编报规则第9号》的规定计算。}

\subsection{管理层讨论与分析}

\subsubsection{资产结构分析}

公司资产总额随经营规模的扩大持续增长，报告期各期末分别为96,384.52万元、
128,475.36万元和168,392.45万元，以流动资产为主。2025年12月31日，货币资金、应
收账款、存货占流动资产的比例分别为25.63\%、29.29\%和23.86\%，资产结构与公司
项目制、定制化的业务模式相匹配。非流动资产的增长主要系算力集群建设投入及研
发设备购置所致。

\subsubsection{偿债能力分析}

报告期各期末，公司流动比率分别为1.93倍、1.87倍和1.84倍，速动比率分别为1.42
倍、1.39倍和1.40倍，短期偿债能力良好；资产负债率分别为39.58\%、41.56\%和
42.36\%，随经营规模扩大及银行借款增加小幅上升，整体处于合理水平。报告期内，
公司不存在到期未偿还的债务，与主要往来银行保持了良好的合作关系。

\subsubsection{经营成果分析}

报告期各期，公司营业收入分别为82,456.32万元、105,238.47万元和132,594.18万
元，2024年度、2025年度分别较上年增长27.63\%和26.00\%，主要受益于下游行业需
求增长及公司产品竞争力提升。同期，公司综合毛利率分别为42.35\%、43.18\%和
44.02\%，稳中有升，主要系毛利率较高的大模型解决方案标杆项目占比提升及开放平
台服务规模效应逐步显现所致。

报告期各期，公司期间费用情况如下：

\biaodw{万元}
\begin{gstab}
\begin{longtable}{|L{46mm}|R{30mm}|R{30mm}|R{30mm}|}
\hline
\thA{项目} & \thB{2025年度} & \thB{2024年度} & \thB{2023年度} \\ \hline
\endfirsthead
\hline
\thA{项目} & \thB{2025年度} & \thB{2024年度} & \thB{2023年度} \\ \hline
\endhead
销售费用 & 9,235.46 & 7,562.34 & 6,124.35 \\ \hline
管理费用 & 6,845.28 & 5,632.45 & 4,652.36 \\ \hline
研发费用 & 12,890.73 & 9,825.64 & 7,412.35 \\ \hline
财务费用 & 285.64 & 196.35 & 152.46 \\ \hline
{\bfseries 期间费用合计} & {\bfseries 29,257.11} & {\bfseries 23,216.78} & {\bfseries 18,341.52} \\ \hline
占营业收入的比例 & 22.07\% & 22.06\% & 22.24\% \\ \hline
\end{longtable}
\end{gstab}

报告期内，公司期间费用率保持稳定，费用增长与业务规模扩张相匹配，其中研发费
用增幅最大，体现了公司对技术研发的持续投入。

\subsubsection{现金流量分析}

报告期各期，公司经营活动产生的现金流量净额分别为10,236.45万元、14,528.36万
元和19,435.82万元，与净利润规模基本匹配，回款质量良好；投资活动现金流出主
要用于智算集群建设、研发设备购置及使用闲置资金进行现金管理；筹资活动现金流入主
要为新增银行借款。

\subsubsection{重大资本性支出及资产负债表日后事项}

报告期内，公司资本性支出主要投向智算集群建设与研发测试设备购置。截至本招股
说明书签署日，公司不存在其他应披露的重大资本性支出安排及资产负债表日后重大
事项。

\subsection{税收优惠情况}

公司及子公司示例智算均为高新技术企业，报告期内按15\%的税率缴纳企业所得税，
并依法享受研发费用加计扣除政策。报告期各期，公司享受的税收优惠占利润总额的
比例未超过20\%，公司经营业绩对税收优惠不存在严重依赖。

\subsection{报告期利润分配情况}

报告期内，公司未进行现金分红，亦未实施送红股或资本公积转增股本。截至2025年
12月31日的滚存未分配利润，按照公司2026年第一次临时股东会决议，由本次发行
完成后的新老股东按持股比例共同享有。

\subsection{财务报告审计截止日后的主要财务信息}

财务报告审计截止日（2025年12月31日）后，公司经营情况正常，2026年1—3月经审
阅的主要财务数据详见本招股说明书``重大事项提示''之``四、财务报告审计截止日
后的主要财务信息及经营状况''。

\subsection{盈利预测}

公司未编制盈利预测报告。
